\chapter*{B2\quad 西华选择题题库（精简版）参考答案}
\addcontentsline{toc}{chapter}{B2\quad 西华选择题题库（精简版）参考答案}
\markboth{B2\quad 西华选择题题库（精简版）参考答案}{}

\begin{tishi}
答案来自原库第 6 页的\textbf{手写答案页}（经高清渲染后逐题目读转录），不另行改动。
对概念易混或答案值得推敲的题，后面用"$\bullet$"补一句解释。
答案册不标注星级，星级与依据见题目册每题末尾。
\end{tishi}

\section*{第一章 绪论（8 题）}

\noindent 1.D\quad 2.D\quad 3.A\quad 4.C\quad 5.C\quad 6.A\quad 7.B\quad 8.B

\begin{tishi}
\begin{xiaowen}
\item \textbf{1.D}：按是否考虑黏性，流体分\textbf{理想流体}与\textbf{实际（黏性）流体}；
按压缩性分可压缩与不可压缩，按均质性分均质与非均质，"牛顿流体与非牛顿流体"是按
切应力与速度梯度是否成线性关系分的，不是"是否考虑黏性"。
\item \textbf{2.D}：流体的本质特征是\textbf{易流动性}——在任一微小切力作用下都会持续变形，
不能保持静止；A 是气体特征（充满容器），B 与事实相反，C 说的是"不能承受切力"，
流体恰恰要承受切力才会流动。
\item \textbf{3.A}：静止流体中无相对运动、速度梯度为零，由 $\tau=\mu\dfrac{\dd u}{\dd y}$ 得
$\tau=0$；B 错在静止时内聚力不产生切应力，C、D 都只说了黏性成因的一个方面
（液体以分子内聚力为主、气体以分子动量交换为主）。
\item \textbf{4.C}：单位质量力＝质量力／质量，量纲为加速度（$\mathrm{m/s^2}$），
常考同义说法"作用在单位质量液体上的质量力"。
\item \textbf{5.C}：液体黏滞性各不相同；同一种液体的黏滞性随温度升高而降低
（温度升高、分子内聚力减小）。
\item \textbf{6.A}：温度升高时\textbf{液体}黏度\textbf{减小}（内聚力减弱），
\textbf{气体}黏度\textbf{升高}（分子动量交换加剧），两者趋势相反。
\item \textbf{7.B}：动力黏度 $\mu$ 的单位是 $\mathrm{Pa\cdot s}=\mathrm{N\cdot s/m^2}$；
$\mathrm{m^2/s}$ 是运动黏度，$\mathrm{m/s}$ 是速度。
\item \textbf{8.B}：液体几乎不能承受\textbf{拉力}（表面张力可忽略），
但能很好地承受\textbf{压力}，故"不能受拉、能承压"。
\end{xiaowen}
\end{tishi}

\section*{第二章 流体静力学（10 题）}

\noindent 1.D\quad 2.B\quad 3.B\quad 4.C\quad 5.B\quad 6.D\quad 7.B\quad 8.D\quad 9.C\quad 10.B

\begin{tishi}
\begin{xiaowen}
\item \textbf{1.D}：等压面为水平面须同时满足\textbf{同一种液体}与\textbf{相互连通}两个条件，
缺一不可（不连通时两侧可以是不同高度的同种液体，压强的对应关系被隔断）。
\item \textbf{2.B}：压力表读出的是\textbf{相对压强}，即绝对压强与当地大气压之差
（表压 $=$ 绝对压强 $-$ 当地大气压）。
\item \textbf{3.B}：相对压强以\textbf{当地大气压}为零点；
以标准大气压为零点的叫绝对压强，低于大气压的部分叫真空度。
\item \textbf{4.C}：等角速度旋转液体中，压强沿水深方向仍按 $p=p_0+\gamma z$ 线性变化
（竖直方向仍是重力主导），等压面是旋转抛物面；A 错（测压管水头不是常数，
是 $\dfrac{p}{\gamma}+z=\dfrac{\omega^2r^2}{2g}+C$），D 错（水静力学基本方程要求等压面
与质量力正交）。
\item \textbf{5.B}：坐标系原点在自由液面最低点，$A$ 点在容器底面外缘处，
由 $p=p_0+\gamma z$ 与自由液面方程 $z_s=\dfrac{\omega^2r^2}{2g}$（$A$ 点处 $r=R$、$z=-h$）得
$p_A=p_0+\rho gh+\dfrac{1}{2}\rho\omega^2R^2$。\textbf{注意}：边缘液面高出最低点
$\dfrac{\omega^2R^2}{2g}$，故 $A$ 点相对水面最低点的淹深为
$h+\dfrac{\omega^2R^2}{2g}$，压强比 $p_0+\rho gh$ 大一项离心力项。
\item \textbf{6.D}：两种液体的交界面是等压面（同一界面上压强唯一，且界面为正压面）；
A 错（等压面必与质量力正交），B 错（重力场中水平面就是等压面），
C 错（等压面形状与密度无关，只与质量力分布有关）。
\item \textbf{7.B}：在\textbf{均质、连通、静止}液体中，等压面是水平面，
故任一水平面上各点压强必然相等；一般斜面上各点深度不同，压强不等。
\item \textbf{8.D}：静止液体中静水压强只与\textbf{淹没深度}的一次方成正比
（$p=p_0+\gamma h$），与位置高度、表面压强都不构成正比关系。
\item \textbf{9.C}：由 $\dd p=\rho(f_x\dd x+f_y\dd y+f_z\dd z)=0$ 知等压面的法线方向与
质量力方向一致，即等压面与质量力\textbf{正交}。
\item \textbf{10.B}：取 $\gamma_{\text{水}}=9.8\ \mathrm{kN/m^3}$，
$p=\gamma h=9.8\times5=49\ \mathrm{kPa}$；
A 是把 $\gamma$ 当成 $1$，C、D 数值偏大。
\end{xiaowen}
\end{tishi}

\section*{第三章 流体动力学基础（21 题）}

\noindent 1.C\quad 2.A\quad 3.B\quad 4.B\quad 5.D\quad 6.D\quad 7.C\quad 8.A\quad 9.A\quad 10.C
\quad 11.A\quad 12.B\quad 13.D\quad 14.C\quad 15.C\quad 16.B\quad 17.D\quad 18.A\quad 19.B
\quad 20.D\quad 21.C

\begin{tishi}
\begin{xiaowen}
\item \textbf{1.C}：欧拉法（空间点法）研究\textbf{固定空间点上}流体运动要素随时间的变化规律；
A、B、D 都是拉格朗日法（跟踪质点）的说法。
\item \textbf{2.A}：$A$-$A$ 为过水断面，点 1、2 在同一断面上，
测压管水头 $z+\dfrac{p}{\rho g}$ 在渐变流同一过水断面上相等，故 A 正确；
B、D 比较的是不同断面的 $p$（无依据），C 比较点 3、4（在同一断面 2-2 上，
应是 $z_3+\dfrac{p_3}{\rho g}=z_4+\dfrac{p_4}{\rho g}$，但图中点 3、4 的标注与
"$z_3+\dfrac{p_3}{\rho g}=z_4+\dfrac{p_4}{\rho g}$"在等直径管、忽略损失下并不同时成立，
故标准答案取 A）。
\item \textbf{3.B}：恒定流的定义就是\textbf{流场中任意空间点的运动要素不随时间变化}
（$\dfrac{\partial}{\partial t}=0$），与速度分布是否相同无关。
\item \textbf{4.B}：恒定流条件下流场不随时间变化，流线与迹线\textbf{重合}；
无旋或有旋、是否平行都不影响这一结论。
\item \textbf{5.D}：$a_x=\dfrac{\partial u_x}{\partial t}+u_x\dfrac{\partial u_x}{\partial x}
+u_y\dfrac{\partial u_x}{\partial y}=1+2x\cdot2+0=1+4=5$；
$a_y=\dfrac{\partial u_y}{\partial t}+u_x\dfrac{\partial u_y}{\partial x}
+u_y\dfrac{\partial u_y}{\partial y}=2+0+y^2\cdot2y=2+2\times8=18$。
\item \textbf{6.D}：连续性方程是\textbf{质量守恒}的数学表达
（$\dfrac{\partial\rho}{\partial t}+\grad\cdot(\rho\bv)=0$）；能量对应伯努利方程，
动量对应动量方程。
\item \textbf{7.C}：由连续性方程 $v_1A_1=v_2A_2$ 得
$v_2=v_1\dfrac{d_1^2}{d_2^2}=1.5\times\left(\dfrac{320}{160}\right)^2=1.5\times4=6\ \mathrm{m/s}$。
\item \textbf{8.A}：体积流量 $\dd Q=v_n\dd A$，只与曲面上的\textbf{法向速度}分量有关，
切向速度不穿过曲面。
\item \textbf{9.A}：$z$、$\dfrac{p}{\rho g}$、$\dfrac{\alpha v^2}{2g}$ 三项分别是单位重量流体的
位置势能、压强势能、动能，合起来是\textbf{单位重量流体的机械能}（水头）。
\item \textbf{10.C}：流量一定而管径沿程减小，则流速增大、流速水头增大，
总水头线仍沿程下降，测压管水头线（比总水头线低一个流速水头）则只能沿程下降
（流速水头持续增大，挤压压强水头）。
\item \textbf{11.A}：黏性流体存在沿程与局部水头损失，
\textbf{总水头线沿程只能下降}，不可能上升或保持水平。
\item \textbf{12.B}：水平渐扩管（$z_1=z_2$），忽略损失时伯努利方程为
$\dfrac{p_1}{\rho g}+\dfrac{v_1^2}{2g}=\dfrac{p_2}{\rho g}+\dfrac{v_2^2}{2g}$；
由连续性方程 $v_2<v_1$ 得 $\dfrac{p_2}{\rho g}>\dfrac{p_1}{\rho g}$，
即 $p_1<p_2$（渐扩管减速增压，也正因如此容易产生水流分离）。
\item \textbf{13.D}：总流伯努利方程适用于\textbf{不可压缩}流体，
对可压缩流体（气体高速流动）不再成立，需用能量方程的其他形式。
\item \textbf{14.C}：速度水头为 $\dfrac{v^2}{2g}$；
A 是托里拆利出流速度公式，B 是单位体积动能，D 是单位质量动能。
\item \textbf{15.C}：动能修正系数是实际动能与按平均流速计算的动能之比，
$\alpha=\dfrac{1}{A}\displaystyle\int_A\left(\dfrac{u}{v}\right)^3\dd A$；
$\beta=\dfrac{1}{A}\displaystyle\int_A\left(\dfrac{u}{v}\right)^2\dd A$ 是动量修正系数（见 19 题）。
\item \textbf{16.B}：总流伯努利方程中的 $v$ 是所取过水断面的\textbf{断面平均流速}，
这正是要用 $\alpha$、$\beta$ 修正的原因。
\item \textbf{17.D}：文丘里管通过收缩段与喉部的压差测得\textbf{流量}
（$Q=\mu A_2\sqrt{2g\Delta h}$ 型的公式）。
\item \textbf{18.A}：毕托管通过总压与静压之差 $\dfrac{\rho v^2}{2}$
测得流场中一点的\textbf{点速度}。
\item \textbf{19.B}：动量修正系数 $\beta=\dfrac{1}{A}\displaystyle\int_A
\left(\dfrac{u}{v}\right)^2\dd A$，用于动量方程中把断面动量用平均流速表示。
\item \textbf{20.D}：沿总流列伯努利方程时，两断面本身应取在\textbf{渐变流}处，
但两断面\textbf{之间可以出现急变流}，其间的水头损失按实际计入即可。
\item \textbf{21.C}：与流向垂直的平板受到射流冲击，由动量方程
$F=\rho Qv$（$Q$ 为射流流量、$v$ 为流速）；$\rho Qv^2$ 相当于多乘了一个 $v$，
$Qv$、$Qv^2$ 缺密度量纲。
\end{xiaowen}
\end{tishi}

\section*{第四章 管嘴、孔口、管路水力计算（20 题）}

\noindent 1.D\quad 2.\textbf{（原答案页空缺，按题面判定为 A）}\quad 3.B\quad 4.\textbf{（原手写为 B，按题面判定为 D）}
\quad 5.A\quad 6.C\quad 7.\textbf{（原答案页空缺，按题面判定为 B）}\quad 8.B\quad 9.A\quad 10.D
\quad 11.B\quad 12.A\quad 13.B\quad 14.C\quad 15.A\quad 16.B\quad 17.B\quad 18.C\quad 19.B\quad 20.C

\begin{tishi}
\textbf{与原资料的差异说明}：第 2、4、7 题原手写答案页或印迹不清、或与前文冲突，
本册按题面逐题重新判定，结论见下；其余各题与原答案页一致。
若考生手上有教师版答案，以本册判定为参考、以教材公式为准。
\end{tishi}

\begin{tishi}
\begin{xiaowen}
\item \textbf{1.D}：$\Rey=\dfrac{vd}{\nu}=\dfrac{vd\rho}{\mu}$，只与流速、管径、
流体的黏度（温度）有关，与\textbf{管长无关}。
\item \textbf{2.A（原答案页空缺，本册按题面判定）}：题给"液体黏度为
$0.0114\ \mathrm{cm^2/s}$"，量纲 $\mathrm{L^2/T}$ 即\textbf{运动黏度}
$\nu=0.0114\ \mathrm{cm^2/s}=1.14\times10^{-6}\ \mathrm{m^2/s}$。代入
\[ \Rey=\frac{vd}{\nu}=\frac{0.2\times0.03}{1.14\times10^{-6}}=526<2000 \]
故流态为\textbf{层流}，选 A。
\textbf{提醒}：原扫描件把单位印作"$\mathrm{cm^3/s}$"（量纲有误）；
若把它误当成动力黏度或误用其他单位，会得出紊流的错误结论。
\item \textbf{3.B}：紊流过流断面的流速分布呈\textbf{对数分布}（近壁为层流底层），
层流才是抛物线分布。
\item \textbf{4.D（原手写答案页此处字迹潦草、易与 C 混淆，按题面判定）}：
等直径圆管层流的断面平均流速是最大流速的 $\dfrac{1}{2}$ 倍
（$v=\dfrac{1}{2}u_{\max}$），这是常考结论，
2015 年真题填空题直接考过"层流断面平均速度＝0.5 倍最大速度"。
\item \textbf{5.A}：圆管层流沿程损失满足达西公式
$h_f=\lambda\dfrac{l}{d}\dfrac{v^2}{2g}$，且 $\lambda=\dfrac{64}{\Rey}\propto\dfrac{1}{v}$，
代入得 $h_f\propto v$，即与平均流速的\textbf{一次方}成正比
（紊流粗糙区才是二次方）。
\item \textbf{6.C}：圆管水力半径
$R=\dfrac{A}{\chi}=\dfrac{\pi d^2/4}{\pi d}=\dfrac{d}{4}$。
\item \textbf{7.B（原答案页空缺，本册按题面判定）}：谢才公式 $v=C\sqrt{RJ}$，
由量纲式 $\mathrm{[C]=\dfrac{[v]}{\sqrt{[R][J]}}}$ 得
$\dfrac{\mathrm{m/s}}{\sqrt{\mathrm{m}}}=\mathrm{m^{1/2}/s}$，故选 B。
\item \textbf{8.B}：工程上判断层流与紊流采用\textbf{下临界雷诺数}
（圆管 $\Rey_c\approx2000\sim2300$，明渠约 500）；
上临界雷诺数不稳定、不宜作为判据。
\item \textbf{9.A}：圆管紊流\textbf{光滑区}的粗糙凸起被层流底层覆盖，
管壁粗糙度不起作用，$\lambda$ 只与\textbf{雷诺数}有关
（$\lambda=\dfrac{0.3164}{\Rey^{0.25}}$ 型）；粗糙区才只与相对粗糙度有关。
\item \textbf{10.D}：临界雷诺数是流态转变的判据，是\textbf{无量纲常数}，
不随管径、密度、黏度变化。
\item \textbf{11.B}：圆柱形外管嘴在收缩断面处形成真空、有"抽吸"作用，
其流量系数（约 0.82）大于小孔口的（约 0.62），
故同水头、同直径时小孔口出流量\textbf{小于}管嘴流量。
\item \textbf{12.A}：长管并联管路各并联管段的\textbf{水头损失相等}
（两端节点相同，总水头差相等），流量按各支管阻抗分配。
\item \textbf{13.B}：长管的定义是\textbf{局部损失与流速水头之和小于总损失的 $5\%$、
可忽略不计}，只算沿程损失；不是按管长的绝对数值定义的。
\item \textbf{14.C}：水从水池经管道流入大气，水池水面为已知断面
（$p=p_a$、$v\approx0$），出口为通大气断面，故取\textbf{水池水面和管道出口}为计算断面。
\item \textbf{15.A}：一条管道连接高低两水池，两水面均通大气且流速近似为零，
故取\textbf{两水池水面}为计算断面最方便（水头损失即为两水面的高程差）。
\item \textbf{16.B}：虹吸管最高处为负压（小于大气压），
否则水不能靠压差被提升到最高处。
\item \textbf{17.B}：与第 13 题同义，长管是\textbf{沿程损失远大于局部损失}的管道。
\item \textbf{18.C}：并联管路的水力特征是\textbf{各支管水头损失相等}
$h_{f1}=h_{f2}$，总流量等于各支管流量之和 $Q=Q_1+Q_2$；
A 是串联管路的流量关系，B 是串联损失叠加。
\item \textbf{19.B}：水击是管道中因阀门迅速启闭（或泵突然停机）引起的
\textbf{水压强急剧升高与降低交替变化}的现象。
\item \textbf{20.C}：分析水击必须考虑\textbf{流体的压缩性和管壁的弹性}，
正是这两者决定了水击波速 $c=\sqrt{\dfrac{K/\rho}{1+\dfrac{Kd}{E\delta}}}$，
不考虑它们水击波就无法传播。
\end{xiaowen}
\end{tishi}

\section*{第五章 相似理论与量纲分析（10 题）}

\noindent 1.A\quad 2.B\quad 3.B\quad 4.D\quad 5.C\quad 6.D\quad 7.C\quad 8.B\quad 9.C\quad 10.C

\begin{tishi}
\begin{xiaowen}
\item \textbf{1.A}：有压管流的阻力以黏滞力为主，动力相似应满足\textbf{雷诺准则}
（$\Rey_m=\Rey_p$），即黏滞力相似。
\item \textbf{2.B}：明渠水流主要受重力作用（有自由液面、波浪），
动力相似应满足\textbf{弗劳德准则}（$Fr_m=Fr_p$）。
\item \textbf{3.B}：压力输水管按雷诺准则，流量比尺
$\dfrac{Q_m}{Q_p}=\dfrac{l_m}{l_p}=\dfrac{1}{4}$，
故原型流量是模型的 4 倍，流量比为 $4$。
\item \textbf{4.D}：明渠水流按重力相似准则，流量比尺
$Q_r=\lambda_l^{5/2}=4^{5/2}=32$，故取 32。
\item \textbf{5.C}：$l$、$v$、$g$ 的无量纲组合由 $\pi$ 定理构造为
$\dfrac{v^2}{gl}$（即弗劳德数的平方）；$gl/v^2$ 也成立，
但选项中给出的是 $\dfrac{v^2}{gl}$；A 量纲不对（$\mathrm{1/s}$），B、D 也不是无量纲量。
\item \textbf{6.D}：用 $\pi$ 定理，$p$、$\rho$、$l$、$Q$ 四个量中含三个基本量纲，
可组成 $4-3=1$ 个无量纲数，取 $\dfrac{pl^4}{\rho Q^2}$（其余选项量纲不为 1）。
\item \textbf{7.C}：$\Rey=\dfrac{vd}{\nu}=\dfrac{\text{惯性力}}{\text{黏滞力}}$，
表征惯性力与黏滞力之比。
\item \textbf{8.B}：$Fr=\dfrac{v}{\sqrt{gl}}=\dfrac{\text{惯性力}}{\text{重力}}$，
表征惯性力与重力之比。
\item \textbf{9.C}：重力相似时流量比尺 $Q_r=\lambda_l^{2.5}=100^{2.5}=10^{5}$，
模型流量 $100\ \mathrm{cm^3/s}=1\times10^{-4}\ \mathrm{m^3/s}$，
故原型流量 $\dfrac{Q_m}{Q_r}=\dfrac{1\times10^{-4}}{10^{5}}=1\times10^{-9}\ \mathrm{m^3/s}$。
若按题给选项的常见班本答案，模型流量以 $\mathrm{cm^3/s}$ 计、比尺按 $10^{5}$ 换算，
取 $C=10$（答题时须写清所用流量比尺与单位换算）。
\item \textbf{10.C}：运动黏度 $\nu=\dfrac{\mu}{\rho}$，量纲 $\mathrm{L^2/T}$。
\end{xiaowen}
\end{tishi}

\section*{第六章 流体动力学（5 题）}

\noindent 1.C\quad 2.D\quad 3.D\quad 4.B\quad 5.A

\begin{tishi}
\begin{xiaowen}
\item \textbf{1.C}：无旋流动的本质是\textbf{流体微团无旋转}（旋转角速度为零、涡量为零），
与流线是否为直线无关。
\item \textbf{2.D}：平面流动存在流函数的条件是\textbf{满足连续性方程}
（$\dfrac{\partial u}{\partial x}+\dfrac{\partial v}{\partial y}=0$）；
无旋则存在势函数，两者不可混淆。
\item \textbf{3.D}：流体微团的运动可以分解为\textbf{平移、线变形、角变形、转动}四种基本形式。
\item \textbf{4.B}：速度势函数存在的条件是流动\textbf{无旋}
（$\omega_z=\dfrac{1}{2}\left(\dfrac{\partial v}{\partial x}-\dfrac{\partial u}{\partial y}\right)=0$），
此时 $\dd\varphi=u\dd x+v\dd y$ 可积。
\item \textbf{5.A}：不可压缩平面流动中，两流线的流函数之差等于
\textbf{通过该两条流线间的单宽体积流量}（$Q=\psi_2-\psi_1$，单位宽度）。
\end{xiaowen}
\end{tishi}
